\chapter{Anger}

\section*{Definition}
Anger is a normal emotion ranging from irritation to intense fury. It may involve physical arousal and can be expressed through words or actions, but feeling angry is not the same as being aggressive or violent. Persistent, disproportionate, or harmful anger may signal a treatable mental, medical, substance-related, or situational problem.

\section*{When to See a Doctor}

\begin{itemize}
 \item Anger results in physical harm to self or others, destruction of property, threats of violence, or fear that control may be lost
 \item Sudden onset of anger without a prior history, especially with confusion or other neurologic symptoms
 \item Accompanying confusion, fever, focal neurologic signs, or recent head trauma
\end{itemize}

\section*{How It Develops}
Perceived threat, frustration, pain, or injustice can activate brain systems involved in threat detection and arousal. Stress, sleep loss, trauma, intoxication or withdrawal, and illness can reduce emotional regulation and increase impulsivity. Anger is shaped by thoughts, learned responses, context, and biology; a simple neurotransmitter explanation cannot predict an individual's behavior.

\section*{Causes}
\begin{itemize}
 \item \textbf{Psychiatric:} Intermittent explosive disorder, depression, bipolar disorder (manic/irritable phase), borderline personality disorder, post-traumatic stress disorder
 \item \textbf{Neurologic:} Traumatic brain injury (especially orbitofrontal cortex), dementia, stroke, epilepsy (temporal lobe), brain tumors
 \item \textbf{Medical:} Hyperthyroidism, hypoglycemia, delirium (any cause), chronic pain, premenstrual dysphoric disorder
 \item \textbf{Substance-related:} Alcohol intoxication or withdrawal, stimulants (cocaine, amphetamines), anabolic steroids, benzodiazepine withdrawal
 \item \textbf{Psychosocial:} Chronic stress, history of abuse, sleep deprivation, interpersonal conflict, perceived injustice
\end{itemize}

\section*{Related Symptoms}
\begin{itemize}
 \item Irritability, agitation, impulsivity, hostility, verbal outbursts
 \item Tachycardia, hypertension, sweating, flushing, muscle tension
 \item Depressive symptoms, anxiety, anhedonia, sleep disturbance
\end{itemize}
\textbf{Important:} Anger is not a DSM-5 diagnosis by itself but rather a symptom of several underlying disorders that must be identified and treated.

\section*{Medical Evaluation}
\begin{itemize}
 \item Thorough history including onset, triggers, duration, frequency, and consequences of anger episodes
 \item Screening for depression, anxiety, PTSD, substance use, and personality disorders
 \item Psychiatric referral for structured assessment and risk evaluation
 \item Targeted laboratory tests such as glucose, thyroid testing, or toxicology may be used when symptoms suggest a medical or substance-related cause; brain imaging is reserved for new neurologic signs, head injury, or other clinical indications
\end{itemize}

\section*{Treatment Options}
\begin{itemize}
 \item Treat underlying psychiatric or medical condition
 \item Cognitive--behavioral therapy (CBT) with anger management: identifying triggers, restructuring thoughts, behavioral rehearsal
 \item Medication is chosen for the underlying condition, not anger alone. SSRIs may treat depression, anxiety, or selected impulse-control disorders; mood stabilizers or antipsychotics are used only for appropriate diagnosed conditions and require clinician monitoring.
 \item Safety planning should include warning signs, coping steps, trusted contacts, emergency services, and reducing access to weapons or other means of harm.
\end{itemize}

\section*{Self-Care and Home Management}
\begin{itemize}
 \item Use de-escalation techniques: remove self from the triggering situation, take deep breaths, count to ten
 \item Practice relaxation strategies: progressive muscle relaxation, guided imagery, mindfulness meditation
 \item Regular sleep, meals, and activity can support emotional regulation; do not exercise or drive when too angry or unsafe
 \item Avoid alcohol and recreational stimulants, and notice whether caffeine worsens symptoms
 \item If there is immediate danger, leave the situation if possible and call emergency services; do not try to handle a violent crisis alone
\end{itemize}

\section*{References}
\begin{thebibliography}{99}
\bibitem{anger:ref1} Siever LJ. ``Neurobiology of aggression and violence.'' \textit{Am J Psychiatry}. 2008;165(4):429-442.
\bibitem{anger:ref2} Coccaro EF. ``Intermittent explosive disorder as a disorder of impulsive aggression.'' \textit{J Clin Psychiatry}. 2012;73(4):479-486.
\bibitem{anger:ref3} Lee R, Coccaro EF. ``The neuropsychopharmacology of criminality and aggression.'' \textit{Can J Psychiatry}. 2013;58(9):508-516.
\bibitem{anger:ref4} Sadock BJ, Sadock VA, et al. \textit{Kaplan \& Sadock\'s Comprehensive Textbook of Psychiatry}, 10th ed. Wolters Kluwer, 2017.
\bibitem{anger:ref5} American Psychiatric Association. \textit{Diagnostic and Statistical Manual of Mental Disorders}, 5th ed. APA, 2013.
\bibitem{anger:ref6} National Health Service. Get help with anger. Updated 2025.
\end{thebibliography}
